% !TeX program = xelatex
\documentclass[UTF8,a4paper,12pt]{ctexart}

% 页面与基础排版
\usepackage{geometry}
\geometry{left=2.8cm,right=2.8cm,top=2.6cm,bottom=2.6cm}
\usepackage{setspace}
\onehalfspacing
\usepackage{indentfirst}
\setlength{\parindent}{2em}
\usepackage{array}
\usepackage{tabularx}
\usepackage{booktabs}
\usepackage{amssymb}
\usepackage{enumitem}
\usepackage{hyperref}
\usepackage{multirow}
\hypersetup{hidelinks}
\hypersetup{hidelinks}
\usepackage{graphicx}
\usepackage{float}
\usepackage{caption}
\usepackage{listings}
\usepackage{xcolor}

% 代码块样式
\definecolor{codebg}{RGB}{245,245,245}
\lstset{
  basicstyle=\ttfamily\small,
  backgroundcolor=\color{codebg},
  frame=single,
  numbers=left,
  numberstyle=\tiny,
  breaklines=true,
  showstringspaces=false,
  tabsize=2
}

% 标题层级格式
\ctexset{
  section = {
    format = \zihao{3}\heiti\centering,
    name = {},
    number = \chinese{section},
    aftername = {、},
    beforeskip = 1.5ex,
    afterskip = 1.2ex
  },
  subsection = {
    format = \zihao{4}\heiti,
    name = {},
    number = \arabic{section}.\arabic{subsection},
    aftername = {\quad}
  },
  subsubsection = {
    format = \zihao{-4}\heiti,
    name = {},
    number = \arabic{section}.\arabic{subsection}.\arabic{subsubsection},
    aftername = {\quad}
  }
}

\newcommand{\placeholder}[1]{\textbf{#1}}
\newcommand{\blankline}{\par\vspace{1.2em}}

\begin{document}

\begin{center}
    {\zihao{2}\heiti 计算机操作系统课程设计详细设计文档}
\end{center}

\vspace{2em}

\noindent\textbf{组长：} \underline{邵汉超} \quad
\textbf{学号：} \underline{2023210932}

\vspace{1em}

\noindent\textbf{组员及分工：}

\vspace{0.5em}

\begin{center}
\begin{tabularx}{0.95\textwidth}{|>{\centering\arraybackslash}p{2.5cm}|>{\centering\arraybackslash}p{3cm}|X|}
\hline
\textbf{姓名} & \textbf{学号} & \multicolumn{1}{c|}{\textbf{分工}} \\
\hline
\multirow{6}{*}{邵汉超} & \multirow{6}{*}{2023210932} & 
\textbf{消息总线与内核路由}：设计 MessageBus trait、LockedBus、Envelope、KernelMsg 枚举（7 变体）、Kernel 路由器（消息分发到 4 个服务 channel）\\
\cline{3-3}
 & & \textbf{进程管理服务}：ProcessService 事件循环、PCB/TCB 进程控制块、SMP 调度器（5 种策略）、进程生命周期（fork/exec/exit/reap）、系统调用分发表（24 个 syscall → 17 个独立 handler）\\
\cline{3-3}
 & & \textbf{同步原语}：信号量 Semaphore（AtomicU32 CAS）、互斥锁 MutexLock（owner 跟踪）、TOCTOU 所有权转移、SyncManager\\
\cline{3-3}
 & & \textbf{内存管理与硬件模拟}：MemoryService（分配/映射/换页）、FrameAllocator、SwapManager；VirtualCPU（20+ 指令、取指解码执行）、MMU（地址转换）、Timer（100Hz）\\
\hline
\multirow{4}{*}{唐博艺} & \multirow{4}{*}{2023210942} & 
\textbf{文件系统与虚拟磁盘}：FileService（消息驱动）、VFS 虚拟文件系统（inode 树 + JSON 持久化）、File 文件抽象（内存缓冲 + 磁盘映射）、FileDescriptorManager（每进程 fd 表）、VirtualDisk（512B 扇区）\\
\cline{3-3}
 & & \textbf{设备服务}：DeviceService（直接订阅总线）、Device trait 与 DeviceRegistry、剪贴板设备（ClipboardGet/Set）、驱动管理 DriverManager\\
\cline{3-3}
 & & \textbf{Shell 与用户界面}：交互式命令行（20+ 命令）、pmon 进程监控器（TUI 实时显示）、minigdb 调试器（单步/寄存器/继续）\\
\cline{3-3}
 & & \textbf{用户态测试程序}：syncdemo.asm、schedemo.asm、cmpdemo.asm 及文件/设备操作 .asm 程序\\
\hline
\end{tabularx}
\end{center}

\newpage

% ═══════════════════════════════════════════════════════════
% 一、总体设计
% ═══════════════════════════════════════════════════════════

\section{总体设计}

\subsection{概述}

\subsubsection{功能描述}

genshin-OS（Chao-OS）是一个基于 Rust 语言从零开发的\heiti{微内核操作系统模拟器}。系统在单个 Linux 进程中完整模拟了操作系统的四层抽象：硬件层（虚拟 CPU、MMU、定时器、磁盘）、服务层（进程管理、内存管理、文件系统、设备驱动）、交换层（消息总线与内核路由）、用户交互层（Shell 与 TUI 监控器）。

\textbf{硬件模拟}：VirtualCPU 实现了 20 余条自定义 Mock ISA 指令（含 MOV/ADD/SUB/MUL/DIV/CMP/JMP/JZ/JNZ/LOAD/STORE/INT/HALT），支持取指-解码-执行三阶段流水线、CPU 标志位（Zero/Sign/Carry/Overflow）以及 INT 0x80 软中断。MMU 管理 per-PID 的页表，实现虚拟地址到物理地址的转换。Timer 以 100Hz 频率通过独立线程发出时钟中断，作为进程调度的"心跳"。VirtualDisk 通过文件 I/O 模拟 512 字节扇区的块设备存储。

\textbf{进程管理}：实现完整的\heiti{Fork-Exec-Exit-Wait} 进程生命周期模型，摒弃简单粗暴的 Spawn 模式。每个用户程序（包括 ls、mkdir、cat 等内置命令）都必须经过 Fork 创建→Exec 加载→运行→Exit 退出→Wait 回收的完整流程。PCB（进程控制块）管理进程状态机（Creating→Ready→Running→Blocked→Zombie），支持多线程 TCB。SMP 双核调度器支持 RoundRobin/FIFO/Priority/SJF/MLFQ 五种策略，per-CPU 量子管理和跨核去重。

\textbf{同步原语}：实现了基于 AtomicU32 CAS 循环的信号量（Semaphore，P/V 操作）和基于 AtomicU64 owner 跟踪的互斥锁（MutexLock，支持递归）。关键的 TOCTOU 所有权转移机制——信号量释放时不盲自 count++，而是扫描进程表将所有权直接转移给阻塞的等待者——解决了多核环境下"释放后立即被调用者抢回"的竞态条件。全局信号量 0 和互斥锁 0 在系统启动时预创建，供 syncdemo 等演示程序使用。

\textbf{文件系统}：实现四层文件系统栈——FileService（消息驱动的顶层服务）→VFS（基于 HashMap 的 inode 树，支持文件和目录的创建/查找/删除，通过 Serde JSON 持久化到 .genshin-vfs.json）→File（Vec<u8> 内存缓冲 + VirtualDisk 扇区映射）→VirtualDisk（基于 .genshin-disk.img 文件的 512B 块设备）。FileDescriptorManager 维护每进程的 fd 表，支持 open/close/read/write 操作。启动时自动从 programs/ 目录导入 .asm 源文件到 VFS。

\textbf{IPC 与设备}：IPCManager 提供进程间消息队列（send/receive/peek）和共享内存（create/attach/detach，引用计数管理）。DeviceService 直接订阅消息总线（不经 Kernel 路由），实现剪贴板设备的读写（ClipboardGet/Set），通过 DeviceRegistry 注册和管理设备驱动。

\textbf{用户界面}：交互式 Shell 支持 20+ 条命令（文件操作：ls/tree/mkdir/touch/cat/write/rm/stat；进程管理：ps/pstree/kill/dual；系统：cd/pwd/verbose/uptime/help/exit）。pmon 提供 TUI 实时进程/内存/磁盘监控。minigdb 提供独立 CPU+MMU 的单步调试器（step/regs/continue/quit），可直接调试 .asm 程序而不干扰运行中的系统。

\textbf{项目规模}：约 15,000 行 Rust 代码，211 次 Git 提交，12 份设计文档，20+ 个 .asm 用户态测试程序。

\subsubsection{运行环境}

\begin{itemize}
    \item \textbf{操作系统}：Linux（在 Arch Linux 上开发，兼容 Ubuntu 22.04+）。依赖 Linux 的线程模型（std::thread）、文件 I/O、以及 ANSI 终端控制序列。
    \item \textbf{Rust 运行时}：项目编译为单一二进制文件，无外部运行时依赖（如 JVM、Python 解释器）。所有依赖均为 Rust 生态内的纯 Rust 库。
    \item \textbf{终端要求}：支持 ANSI 转义序列的终端模拟器（如 Alacritty、Kitty、GNOME Terminal、Windows Terminal）。pmon TUI 使用 \texttt{\textbackslash{}033[} 控制码实现光标定位和颜色。
    \item \textbf{外部依赖（Cargo.toml）}：crossbeam-channel（多生产者多消费者无界通道，用于消息总线）、serde/serde\_json（VFS 树的 JSON 序列化/反序列化）。均为 MIT/Apache-2.0 双重许可的工业级库。
    \item \textbf{内存占用}：运行时约 50-100MB（含 4MB 物理内存模拟、两个 4096 页表、VFS 树、磁盘缓存）。
\end{itemize}

\subsubsection{开发环境}

\begin{itemize}
    \item \textbf{编程语言}：Rust（edition 2021）。选择 Rust 的核心原因：① 所有权系统和借用检查器在编译期消除数据竞态和悬垂指针，适合并发密集的操作系统开发；② 零成本抽象的 enum/match/Arc<Mutex<T>> 模式完美匹配消息驱动的架构；③ Cargo 的依赖管理和测试框架无需 Makefile/CMake 等额外构建工具。
    \item \textbf{编译器}：rustc 1.70+（stable channel）。使用 2021 edition 的闭包捕获规则和 const 泛型等特性。
    \item \textbf{构建系统}：Cargo。常用命令：\texttt{cargo build}（调试构建）、\texttt{cargo run}（构建并运行）、\texttt{cargo check}（仅类型检查不编译）、\texttt{cargo test}（运行单元测试）、\texttt{cargo fmt}（代码格式化）、\texttt{cargo clippy}（静态分析 lint）。
    \item \textbf{版本控制}：Git。采用 conventional commits 规范（\texttt{fix:/feat:/docs:/refactor:/chore:}），每个逻辑修改独立提交。仓库共 211 次提交，提交信息包含问题分析和修复说明。
    \item \textbf{编辑器}：VS Code（rust-analyzer 插件提供实时类型检查和代码补全）或 Neovim（rust-tools.nvim）。
    \item \textbf{文档编译}：LaTeX 报告使用 xelatex 编译（ctexart 文档类），需安装 TeX Live 完整版或 texlive-xetex + texlive-latex-extra 包。
\end{itemize}

\subsection{设计思想}

\subsubsection{软件设计构思}

\textbf{1. 微内核架构：最小化信任计算基础}

本项目严格遵循\heiti{微内核（Microkernel）}设计哲学：内核（Kernel）仅做一件事——消息路由。系统调用、进程调度、内存分配、文件操作、设备驱动均以独立的\heiti{用户态服务}形式运行在各自的线程中。这种设计的核心优势是\heiti{隔离性}：一个服务的崩溃（如 FileService 在解析路径时 panic）不会影响其他服务（ProcessService 仍可调度进程）。L4 家族（seL4、Pistachio）和 QNX 等工业级微内核均采用此架构。

内核路由器的实现极为精简（仅 51 行代码，见 src/services/kernel.rs）：它在消息总线上阻塞等待 Envelope，根据 KernelMsg 枚举的变体将消息转发到四个专用通道（process\_tx、intr\_tx、memory\_tx、file\_tx）。内核不解析消息内容、不维护任何状态、不执行业务逻辑——它是真正的\heiti{"消息交换机"}。

\textbf{2. 消息总线：统一通信语言}

所有跨模块通信强制通过 MessageBus trait 定义的接口：send()（发后即忘）、send\_request()（请求-响应）、subscribe()（订阅）。LockedBus 实现基于 Arc<Mutex<Vec<Sender<Envelope>>>>——订阅者列表由 Mutex 保护，消息通过 crossbeam-channel 的 unbounded 通道传递。

KernelMsg 枚举（7 个变体）是整个系统的\heiti{通信语言}：Syscall（用户系统调用）、Interrupt（硬件中断，含 Timer/PageFault/SyscallTrap）、Process（进程服务请求，含 30+ 子变体覆盖调度/IPC/同步/生命周期）、Memory（内存分配/映射/换页）、File（文件 CRUD/目录操作）、Device（设备读写）。任何模块只能使用这个枚举中的类型发送消息，不能调用其他模块的函数。这确保了\heiti{接口的强制统一性}——审查一个模块的行为只需检查它如何处理 KernelMsg。

Envelope 信封机制支持两种通信模式：fire\_and\_forget 用于通知（如 Timer 发送时钟中断），with\_response 用于查询（如 Shell 的 send\_and\_wait 发送 ps 请求并等待进程列表）。响应通过 Envelope 内嵌的 response\_channel（一个 crossbeam Sender<Response>）传递，调用方用 recv\_timeout(3s) 阻塞等待。

\textbf{3. 两层译码：硬件与内核的职责分离}

系统调用的处理采用\heiti{两层译码}模型：\begin{enumerate}[label=(\arabic*)]\item \textbf{CPU 译码指令}：VirtualCPU 执行 INT 0x80 时，仅保存寄存器到 syscall\_regs、设置 syscall\_pending=true、停机。它不知道 R0=18 代表 listdir，也不关心。\item \textbf{ProcessService 译码调用号}：定时器中断的 step loop 检查 syscall\_pending 后，调 handle\_file\_syscall(cpu, r0, r1, r2)，由 match r0 分派到 24 个 syscall\_xxx() handler。\end{enumerate}这种分离的核心理由是\heiti{硬件层不应知道操作系统策略}。真实的 x86 CPU 上，INT 0x80 只触发中断切换到内核态，系统调用号的解释完全由内核完成。我们的模拟保持了同样的抽象边界。

\textbf{4. Fork-Exec 模型：为什么不用 Spawn}

进程创建采用 Unix 经典的\heiti{Fork-Exec-Wait}三步模型，而非 Windows 式的 Spawn（CreateProcess 一步到位）。具体流程：ForkProcess（克隆父进程的页表和 CPU 状态，子进程 R0=0 标识其为 child，从父进程 PC 继续执行）→ ExecProcess（替换子进程的页表，加载 .asm 程序，重置 CPU PC=0，加入就绪队列）→ WaitChild（父进程阻塞直到子进程退出为 Zombie，然后回收资源）。

选择 Fork-Exec 而非 Spawn 的设计理由：① \heiti{正交性}——Fork 只做"复制"，Exec 只做"替换"，职责单一，便于测试和调试；② \heiti{灵活性}——Fork 后 Exec 前可以执行任意操作（如重定向文件描述符），为未来的管道和重定向功能预留空间；③ \heiti{教学价值}——这是 Unix 操作系统的核心设计，理解 Fork-Exec 是理解 Unix 进程模型的关键。fork\_impl 创建子进程后\heiti{不立即调度}，等待 exec\_impl 加载代码后才加入就绪队列，避免无代码进程被调度导致的页错误→Zombie 竞态。

\textbf{5. 系统调用的两路分发：何时走总线、何时不走}

系统调用处理分为两条路径：\begin{itemize}\item \textbf{文件/设备类（R0=10-18, 208-211）}：走消息总线。handler 从进程 0x100 读参数，发 KernelMsg::File/Device，Kernel 路由到 FileService/DeviceService，响应返回后写结果到进程 0x200，设 R2=字节数。\item \textbf{进程/同步类（R0=0, 100-102, 200-205）}：不走总线，直接操作本地数据（process\_table、MMU、SyncManager）。\end{itemize}区分依据是\heiti{数据局部性}：进程和同步的数据在 ProcessService 内部，走总线是浪费；文件数据在 FileService，必须走总线。

\textbf{6. 同步与异步的权衡}

当前 INT 0x80 系统调用都是\heiti{同步}的——CPU 停机等待 handler 返回后才继续。这简化了 .asm 程序（线性执行，无需回调），代价是 CPU 在等 I/O 时完全阻塞。真正的微内核（如 seL4）使用异步 IPC 加 Notification 通知机制。鉴于 .asm 程序仅 3-15 条指令、数据量极小（<4KB），同步等待的开销（recv\_timeout(10ms)）可忽略。未来若支持网络栈或大文件，异步模型将更有价值。

\textbf{7. VFS 持久化：JSON + 二进制混合策略}

VFS 树结构使用 Serde JSON 持久化到 .genshin-vfs.json。选择 JSON 的理由：① \heiti{可调试性}——开发者可直接用文本编辑器查看文件系统状态；② \heiti{简单性}——derive 宏零代码实现序列化。代价是比二进制大约 3-5 倍，但 inode 数不超过 200 的小规模文件系统可忽略。磁盘镜像 .genshin-disk.img 使用\heiti{二进制格式}（512B 扇区），因为文件内容本身就是二进制。这种"元数据 JSON + 数据二进制"的混合策略平衡了可调试性和效率。
\subsubsection{关键技术与算法}

\begin{enumerate}[label=(\arabic*)]
    \item \textbf{SMP 多核调度}：双 CPU 独立调度，RR 时间片轮转算法（quantum=3 ticks），dequeue\_next 跳过已在其他 CPU 运行的 PID，每 CPU 维护 last\_running 状态
    \item \textbf{TOCTOU 所有权转移}：信号量释放时不 count++，而是扫描进程表找到 Blocked 等待者，直接转移所有权。避免多核环境下调用者立即抢回信号量的竞态
    \item \textbf{CAS 无锁同步}：信号量和互斥锁的 wait/signal 操作使用 AtomicU32/U64 的 compare\_exchange\_weak 实现，无需互斥锁保护
    \item \textbf{VFS inode 树}：基于 HashMap 的 inode 表 + 父目录 children HashMap 实现路径解析和文件查找，支持 JSON 持久化
    \item \textbf{系统调用直接路径}：INT 0x80 触发的中断不经过消息总线，直接由 ProcessService 的 handle\_file\_syscall 在定时器中断步进循环中处理，避免 SyscallTrap 消息泛滥
\end{enumerate}

\subsubsection{基本数据结构}

全局核心数据结构包括：

\begin{itemize}
    \item \textbf{KernelMsg}：所有消息的统一枚举，包含 Syscall、Interrupt、Process、Memory、File、Device 七个变体，是整个系统的通信语言
    \item \textbf{PCB（Process Control Block）}：进程控制块，记录进程状态（Creating→Ready→Running→Blocked→Zombie）、页表、CPU 状态、父进程、打开文件等
    \item \textbf{VFSNode}：虚拟文件系统节点，包含 inode、类型、权限、父子关系、数据块列表
    \item \textbf{Semaphore / MutexLock}：基于原子变量的同步原语，Semaphore 用 AtomicU32 计数，MutexLock 用 AtomicU64 跟踪持有者
    \item \textbf{Envelope}：消息信封，封装 KernelMsg + 请求 ID + 响应通道，支持 fire-and-forget 和 request-response 两种通信模式
\end{itemize}

\subsection{基本处理流程}

系统的核心处理流程是\heiti{时钟中断驱动的调度循环}：

\begin{enumerate}[label=(\arabic*)]
    \item Timer 以 100Hz 频率向消息总线发送 Interrupt::Timer
    \item Kernel 路由器将中断转发到 ProcessService 的 intr\_rx 通道
    \item ProcessService 的主循环收到中断后调 handle\_timer\_interrupt()
    \item 遍历两个虚拟 CPU：每个 CPU 先调 schedule(cpu\_id) 从就绪队列选进程，再执行最多 3 条指令
    \item 指令执行过程中若遇到 INT 0x80，直接调用 handle\_file\_syscall 分发系统调用
    \item 检查僵尸进程 → 回收 → 循环
\end{enumerate}

文件系统操作流程：用户进程执行 INT 0x80（R0=18 listdir）→ ProcessService 转发 FileRequest::ListDir 到消息总线 → Kernel 路由到 FileService → FileService 调 VFS.lookup\_path 解析路径 → 收集子节点名称 → 通过响应通道返回 → ProcessService 将结果写入进程虚拟内存 0x200 → 进程通过 INT 0x80（R0=2）打印结果。

\subsection{软件的体系结构设计}

\subsubsection{软件体系结构框图}

\begin{center}
\fbox{\begin{minipage}{0.9\textwidth}
\begin{center}
\textbf{genshin-OS 微内核架构}
\end{center}

\vspace{0.5em}

\begin{center}
\begin{tabular}{|c|}
\hline
\textbf{用户交互层（UI）} \\
Shell（命令解析、pmon 监控、minigdb 调试器） \\
\hline
\textbf{交换层（Exchange Layer）} \\
MessageBus（LockedBus）- KernelMsg 枚举 - Envelope \\
\hline
\textbf{服务层（Service Layer）} \\
ProcessService（进程/调度/IPC/同步） \quad MemoryService（页帧/换页） \\
FileService（VFS/文件/fd） \quad DeviceService（剪贴板） \\
\hline
\textbf{硬件层（Hardware Layer）} \\
VirtualCPU（20+ 指令） \quad MMU（页表） \quad Timer（100Hz） \quad VirtualDisk（512B） \\
\hline
\end{tabular}
\end{center}

\vspace{0.5em}
\begin{center}
\small 所有跨层通信均通过 MessageBus，无直接函数调用
\end{center}
\end{minipage}}
\end{center}

\subsubsection{软件主要模块及其依赖关系说明}

\begin{itemize}
    \item \textbf{Kernel}：唯一订阅 MessageBus 的组件。根据 KernelMsg 类型将消息路由到 ProcessService、MemoryService、FileService 的专用 channel。DeviceService 直接订阅总线，不经 Kernel 路由。
    \item \textbf{ProcessService}：最复杂的服务，依赖 SyncManager（同步原语）、IPCManager（消息队列/共享内存）、Scheduler（就绪队列管理）、MMU（地址转换）。接收 ProcessRequest 和 Syscall 两类消息。
    \item \textbf{FileService}：依赖 VFS（目录树）、FileDescriptorManager（fd 管理）、VirtualDisk（块存储）。通过 VFS 的 JSON 序列化实现持久化。
    \item \textbf{MemoryService}：管理 PhysicalMemory，依赖 MMU 完成页表操作和地址映射。
    \item \textbf{DeviceService}：独立订阅总线，管理剪贴板等设备，通过 DeviceRegistry 注册设备驱动。
    \item \textbf{Shell}：运行在主线程，通过 UIContext.send\_and\_wait 与消息总线交互，不依赖任何服务。
\end{itemize}

\subsection{软件数据结构设计}

\subsubsection{全局数据结构说明}

系统中最关键的全局数据结构是 \textbf{KernelMsg 枚举}（定义在 src/messaging/msg.rs:27），它是所有模块间通信的唯一载体。定义如下：

\begin{lstlisting}[]
pub enum KernelMsg {
    Syscall(Syscall),         // 用户系统调用（CreateProcess, ExitProcess 等）
    Interrupt(Interrupt),     // 硬件中断（Timer, PageFault, SyscallTrap 等）
    Process(ProcessRequest),  // 进程服务请求（调度/IPC/同步/生命周期）
    Memory(MemoryRequest),    // 内存服务请求（分配/映射/换页）
    File(FileRequest),        // 文件服务请求（CRUD/目录操作）
    Device(DeviceRequest),    // 设备服务请求（ClipboardGet/Set）
}
\end{lstlisting}

ProcessRequest 枚举包含 30+ 变体，覆盖进程调度（Schedule/Block/Unblock）、IPC（SendMessage/ReceiveMessage/CreateSharedMemory）、同步（CreateSemaphore/WaitSemaphore/AcquireLock）、生命周期（ForkProcess/ExecProcess/WaitChild/Kill）等所有进程相关操作。

\subsubsection{数据结构与系统单元的关系}

\begin{center}
\renewcommand{\arraystretch}{1.35}
\begin{tabularx}{0.95\textwidth}{|>{\centering\arraybackslash}p{3cm}|>{\centering\arraybackslash}X|>{\centering\arraybackslash}X|>{\centering\arraybackslash}X|>{\centering\arraybackslash}X|}
\hline
 & \textbf{ProcessService} & \textbf{FileService} & \textbf{MemoryService} & \textbf{Shell}\\
\hline
PCB（进程控制块） & $\checkmark$ & & & \\
\hline
VFSNode（文件节点） & & $\checkmark$ & & \\
\hline
Semaphore/MutexLock & $\checkmark$ & & & \\
\hline
PageTable（页表） & $\checkmark$ & & $\checkmark$ & \\
\hline
FileDescriptor & & $\checkmark$ & & \\
\hline
KernelMsg（消息） & $\checkmark$ & $\checkmark$ & $\checkmark$ & $\checkmark$ \\
\hline
VirtualDisk（磁盘） & & $\checkmark$ & $\checkmark$ & \\
\hline
\end{tabularx}
\end{center}

\subsection{软件接口设计}

\begin{enumerate}[label=(\arabic*)]
    \item \textbf{操作界面（命令接口）}
    
    用户通过交互式 Shell 输入命令。Shell 解析命令后通过消息总线发送 ProcessRequest 或 FileRequest，等待服务响应（3 秒超时）。支持的命令包括：ls、mkdir、touch、cat、write、rm、tree、stat（文件操作）、ps/pstree（进程管理）、kill（信号）、dual（多进程演示）、pmon（监控）、minigdb（调试）、cd/pwd（导航）、verbose（调试开关）、uptime（时钟计数）、help。

    \item \textbf{系统功能服务界面（程序接口）}
    
    用户进程通过 INT 0x80 软中断请求系统服务。R0 寄存器存放系统调用号，R1/R2 存放参数。ProcessService 的 handle\_file\_syscall 函数根据 R0 值分派到 24 个独立 handler。文件类系统调用（R0=10-18）将结果写入进程虚拟内存 0x200，进程通过 R0=2（print\_str）打印结果。
\end{enumerate}

\subsubsection{外部接口}

\textbf{输入}：键盘输入命令，通过 stdin 读取；程序文件从 programs/ 目录加载 .asm 源文件，由汇编器（src/services/process/assembler.rs）转换为字节码。

\textbf{输出}：屏幕显示通过 stdout 输出 Shell 提示符、命令执行结果、进程输出（PRINT）、调试信息。文件数据持久化到 .genshin-vfs.json（VFS 结构）和 .genshin-disk.img（虚拟磁盘内容）。

\subsubsection{内部接口}

\textbf{消息总线}：所有服务间通信的唯一通道。发送方调用 bus.send(msg) 或 bus.send\_request(msg)，消息通过 Kernel 路由到目标服务的接收 channel。请求-响应模式下，Envelope 携带 response\_channel，服务通过 envelope.respond\_success() 回复。

\textbf{进程虚拟内存约定}：执行 exec 时，程序路径写入虚拟地址 0x100，参数写入 0x200。文件系统调用的返回数据写入 0x200，字节数写入 CPU 的 R2 寄存器。

\newpage

% ═══════════════════════════════════════════════════════════
% 二、用户界面设计
% ═══════════════════════════════════════════════════════════

\section{用户界面设计}

\subsection{界面的关系图}

系统提供两种用户界面：

\begin{enumerate}[label=(\arabic*)]
    \item \textbf{交互式 Shell}：启动后进入命令行界面，显示欢迎信息和提示符 \texttt{genshin-os:/>}。用户输入命令后，Shell 解析并通过消息总线发送请求，同步等待响应后打印结果。
    \item \textbf{pmon 进程监控器}：输入 \texttt{pmon} 命令进入 TUI 界面，实时显示进程状态表、内存使用、磁盘信息。按 \texttt{q} 退出。
\end{enumerate}

\subsection{界面说明}

\subsubsection{Shell 命令行界面}

启动系统后，终端显示：

\begin{lstlisting}
Initializing Genshin-OS microkernel simulation...
✓ Hardware + Message bus
✓ Kernel
✓ Timer (hardware, 100 Hz)
✓ Process service
✓ Memory service
✓ File service
✓ Device service
✓ Shell

╔════════════════════════════════════════════════════════════╗
║           Welcome to Genshin-OS Microkernel Shell          ║
║                     Version 0.2.0                          ║
╠════════════════════════════════════════════════════════════╣
║  Type 'help' for available commands or 'exit' to quit.     ║
╚════════════════════════════════════════════════════════════╝

genshin-os:/>
\end{lstlisting}

常用命令示例：

\begin{itemize}
    \item \texttt{ls} — 列出当前目录
    \item \texttt{mkdir testdir} — 创建目录
    \item \texttt{touch testdir/a.txt} — 创建文件
    \item \texttt{write testdir/a.txt hello} — 写入文件
    \item \texttt{cat testdir/a.txt} — 读取文件
    \item \texttt{ps} — 查看进程树（状态、PID、名称）
    \item \texttt{dual syncdemo} — 启动两个互斥演示进程
    \item \texttt{verbose on} — 开启详细调试输出
    \item \texttt{pmon} — 打开进程监控器
    \item \texttt{minigdb ls} — 单步调试 ls 程序
\end{itemize}

\newpage

% ═══════════════════════════════════════════════════════════
% 三、相关处理流程（程序单元设计）
% ═══════════════════════════════════════════════════════════

\section{相关处理流程}

\subsection{程序单元1：消息总线与内核路由}

\subsubsection{程序单元说明}

消息总线是整个系统的通信中枢。本单元实现 MessageBus trait 及其 LockedBus 实现、Envelope 信封机制、KernelMsg 统一消息枚举、以及 Kernel 路由器。所有跨模块通信均通过消息总线，无直接函数调用。

\subsubsection{数据结构说明}

\begin{itemize}
    \item \textbf{MessageBus trait}（src/messaging/bus.rs:96）：定义 send（发后即忘）、send\_request（请求-响应）、subscribe（订阅）三个核心方法
    \item \textbf{LockedBus}（src/messaging/bus.rs:159）：基于 Arc<Mutex<Vec<Sender<Envelope>>>> 的实现，send 遍历订阅者广播，subscribe 创建 unbounded channel
    \item \textbf{Envelope}（src/messaging/bus.rs:31）：封装 message + request\_id + response\_channel，支持 fire\_and\_forget 和 with\_response 构造
    \item \textbf{KernelMsg}（src/messaging/msg.rs:27）：7 变体统一消息枚举
    \item \textbf{Kernel}（src/services/kernel.rs）：持有 4 个 Sender（process/intr/memory/file），按消息类型路由
\end{itemize}

\subsubsection{算法及流程}

消息发送流程（以 Shell 执行 ls 命令为例）：

\begin{enumerate}[label=(\arabic*)]
    \item Shell 调用 \texttt{context.send\_request(KernelMsg::Process(ForkProcess\{parent\_pid:1\}))}
    \item LockedBus::send\_request 创建 Envelope（含 response\_channel），遍历 subscribers 用 try\_send 广播
    \item Kernel 的 receiver.recv() 收到消息，route() 判断消息类型为 Process → process\_tx.send(envelope)
    \item ProcessService 的 receiver.try\_recv() 收到消息，handle\_envelope 处理 ForkProcess
    \item ProcessService 通过 envelope.respond\_success(Pid(child\_pid)) 发送响应
    \item Shell 的 rx.recv\_timeout(3s) 收到响应 → 继续发送 ExecProcess
\end{enumerate}

\subsubsection{数据存储说明}

无需数据存储。消息在内存中的 unbounded channel 传输，不持久化。

\subsubsection{源程序文件说明}

\begin{itemize}
    \item \texttt{src/messaging/bus.rs} — MessageBus trait、LockedBus、DirectBus、Envelope
    \item \texttt{src/messaging/msg.rs} — KernelMsg、ProcessRequest、FileRequest、IPCMessage 等枚举
    \item \texttt{src/messaging/response.rs} — Response、ResponseData、ServiceError
    \item \texttt{src/messaging/mod.rs} — 模块导出
    \item \texttt{src/services/kernel.rs} — Kernel 路由器
\end{itemize}

\subsubsection{函数说明}

\begin{itemize}
    \item \texttt{LockedBus::send(msg)} — 遍历所有订阅者，try\_send 广播消息，fire-and-forget
    \item \texttt{LockedBus::send\_request(msg)} — 创建 Envelope + response\_channel，广播，返回 Receiver<Response>
    \item \texttt{LockedBus::subscribe()} — 创建 unbounded channel，将 Sender 加入 subscribers，返回 Receiver
    \item \texttt{Envelope::fire\_and\_forget(msg)} — 创建无响应通道的信封
    \item \texttt{Envelope::with\_response(msg)} — 创建含响应通道的信封
    \item \texttt{Kernel::route(envelope)} — 根据 KernelMsg 变体将消息发送到对应服务的 Sender
\end{itemize}

\subsection{程序单元2：进程管理服务与调度器}

\subsubsection{程序单元说明}

ProcessService 是系统中最核心的服务，负责进程的完整生命周期管理（创建、加载、运行、阻塞、退出、回收）、SMP 双核调度、IPC 通信、同步原语操作、以及系统调用分发。

\subsubsection{数据结构说明}

\begin{itemize}
    \item \textbf{PCB}（src/services/process/pcb.rs:261）：进程控制块，包含 pid、name、state（状态机）、parent\_pid、priority、threads
    \item \textbf{TCB}（src/services/process/pcb.rs:108）：线程控制块，包含 tid、entry\_point、stack\_pointer、state
    \item \textbf{Scheduler}（src/services/process/scheduler.rs:45）：就绪队列 VecDeque<ReadyQueueEntry>，per-CPU cpu\_current 和 cpu\_ticks，支持 RR/FIFO/Priority/SJF/MLFQ
    \item \textbf{IPCManager}（src/services/process/ipc.rs:240）：消息队列 HashMap<Pid, MessageQueue>，共享内存 HashMap<shmid, SharedMemoryRegion>
    \item \textbf{SyncManager}（src/services/process/sync.rs:389）：信号量和互斥锁的集中管理器，预创建全局 semaphore 0 和 mutex 0
\end{itemize}

\subsubsection{算法及流程}

\textbf{SMP 调度算法}（handle\_timer\_interrupt，service.rs:571-740）：

\begin{enumerate}[label=(\arabic*)]
    \item 每 tick 遍历两个 CPU（cpu\_id = 0, 1）
    \item schedule(cpu\_id)：检查当前进程量子是否用完（3 ticks），用完则重新入队；从就绪队列取下一个，跳过已在其他 CPU 运行的 PID
    \item 状态机：将上次运行的进程从 Running 切换为 Ready，将新进程标记为 Running
    \item CPU 步进循环：最多执行 3 条指令，遇 INT 0x80 直接调用 handle\_file\_syscall
    \item Zombie 检查：若 CPU 停机的进程并非 Blocked，则标记为 Zombie，释放信号量
    \item Reaper：每 tick 扫描一个 Zombie 进程并回收（释放内存、移除 PCB）
\end{enumerate}

\textbf{进程创建流程}（fork+exec，非 Spawn）：

\begin{enumerate}[label=(\arabic*)]
    \item Shell 发送 ForkProcess\{parent\_pid: 1\} → fork\_impl 克隆父进程页表 → 创建子进程 PCB/CPU → \textbf{不立即调度}（等待 exec 加载代码）
    \item Shell 发送 ExecProcess\{pid, executable\} → exec\_impl 加载 .asm 程序 → 写虚拟内存 → 重置 CPU PC=0 → 加入就绪队列
    \item 这种方式避免了子进程在代码加载前被调度执行导致页错误/僵尸的问题
\end{enumerate}

\subsubsection{数据存储说明}

进程控制块和调度器状态均为内存数据结构，不持久化。进程退出时页帧被回收。

\subsubsection{源程序文件说明}

\begin{itemize}
    \item \texttt{src/services/process/service.rs} — ProcessService 主逻辑（事件循环、系统调用分发表、定时器中断处理）
    \item \texttt{src/services/process/pcb.rs} — PCB/TCB 进程控制块
    \item \texttt{src/services/process/scheduler.rs} — SMP 调度器
    \item \texttt{src/services/process/ipc.rs} — IPCManager、MessageQueue、SharedMemoryRegion
    \item \texttt{src/services/process/sync.rs} — Semaphore、MutexLock、SyncManager
    \item \texttt{src/services/process/assembler.rs} — .asm 汇编器
\end{itemize}

\subsubsection{函数说明}

\begin{itemize}
    \item \texttt{handle\_timer\_interrupt()} — 定时器中断处理，SMP 调度 + CPU 步进 + Zombie 回收（service.rs:571-740）
    \item \texttt{fork\_impl(parent\_pid)} — 克隆父进程的内存和 CPU 状态（service.rs:973-1050）
    \item \texttt{exec\_impl(pid, executable, args)} — 加载程序、替换页表、重置 CPU（service.rs:1100-1151）
    \item \texttt{handle\_file\_syscall(cpu, r0, r1, r2)} — 系统调用分发表，17 行 match 分发到 24 个独立 handler（service.rs:1450-1486）
    \item \texttt{Scheduler::schedule(cpu\_id)} — SMP 调度决策（scheduler.rs:127-142）
    \item \texttt{Semaphore::wait()} — CAS 循环实现 P 操作（sync.rs:82-107）
    \item \texttt{MutexLock::try\_acquire(pid)} — CAS 实现获取锁（sync.rs:276-306）
\end{itemize}

\subsection{程序单元3：同步原语}

\subsubsection{程序单元说明}

本单元实现了两个核心同步原语：信号量（Semaphore）和互斥锁（MutexLock）。两者均基于 Rust 的 AtomicU32/AtomicU64 实现无锁并发操作。关键的 TOCTOU 所有权转移机制解决了多核环境下信号量释放后立即被调用者抢回的竞态问题。

\subsubsection{数据结构说明}

\begin{itemize}
    \item \textbf{Semaphore}（sync.rs:25）：value（AtomicU32 计数器）、max\_value、wait\_count、owner\_pid、valid 标志
    \item \textbf{MutexLock}（sync.rs:218）：owner（AtomicU64，u64::MAX=未锁定）、count（递归计数）、recursive 标志、wait\_count
    \item \textbf{SyncManager}（sync.rs:389）：HashMap<id, Semaphore> + HashMap<id, MutexLock>，集中管理所有同步原语
\end{itemize}

\subsubsection{算法及流程}

\textbf{Semaphore::wait() P 操作}（sync.rs:82-107）：

\begin{lstlisting}[]
pub fn wait(&self) -> SemaphoreResult {
    let mut current = self.value.load(Ordering::Acquire);
    loop {
        if current == 0 {
            self.wait_count.fetch_add(1, Ordering::AcqRel);
            return SemaphoreResult::WouldBlock;
        }
        match self.value.compare_exchange_weak(
            current, current - 1, Ordering::AcqRel, Ordering::Acquire
        ) {
            Ok(_) => return SemaphoreResult::Acquired,
            Err(new) => current = new,
        }
    }
}
\end{lstlisting}

\textbf{TOCTOU 所有权转移}（service.rs:1668-1693）：
当进程调用 sem\_signal 时，不直接 count++，而是扫描进程表查找 Blocked(WaitingForLock\{lock\_addr == sem\_id\}) 的等待者。找到等待者时：
\begin{enumerate}[label=(\arabic*)]
    \item PCB 状态设为 Ready
    \item scheduler.ready(wpid) 加入就绪队列
    \item cpu.halted = false 唤醒 CPU
    \item 信号量 count 保持 0（不增加）——所有权直接转移给等待者
\end{enumerate}
如果未找到等待者，则 sem.signal() → count++。

\subsubsection{源程序文件说明}

\begin{itemize}
    \item \texttt{src/services/process/sync.rs} — 全文（Semaphore、MutexLock、SyncManager）
    \item \texttt{src/services/process/service.rs:1642-1752} — 同步系统调用 handler
\end{itemize}

\subsection{程序单元4：文件系统与虚拟磁盘}

\subsubsection{程序单元说明}

FileService 是文件系统的顶层服务，整合了 VFS 目录树、File 文件抽象、FileDescriptorManager 文件描述符管理、VirtualDisk 虚拟磁盘四个层次。所有文件操作通过消息总线接收 FileRequest，由 FileService 协调下层组件完成。

\subsubsection{数据结构说明}

\begin{itemize}
    \item \textbf{VFSNode}（vfs.rs:24）：inode、node\_type（File/Directory/SymLink）、name、parent、permissions、size、blocks（磁盘块号列表）、children（子节点 HashMap）
    \item \textbf{VirtualFileSystem}（vfs.rs:173）：nodes: HashMap<inode, VFSNode>，next\_inode: 自增分配器，mounts: 挂载点表
    \item \textbf{File}（file.rs:100）：data: Vec<u8> 内存缓冲、position: 读写光标、start\_sector: 磁盘起始位置、dirty: 脏标记
    \item \textbf{OpenFile}（descriptor.rs:18）：fd + owner + file 引用 + flags
    \item \textbf{FileDescriptorManager}（descriptor.rs:150+）：每进程 fd 表，alloc/dealloc 管理
    \item \textbf{VirtualDisk}（disk.rs:30+）：512B 扇区，read\_sector/write\_sector，基于 .genshin-disk.img 文件
\end{itemize}

\subsubsection{算法及流程}

\textbf{文件读取完整流程}（cat /foo/bar.txt）：

\begin{enumerate}[label=(\arabic*)]
    \item cat.asm：MOV R0,\#10 → INT 0x80（open，路径从 0x100 读取）
    \item ProcessService → bus.send(FileRequest::Open\{path, flags\}) → FileService
    \item FileService：vfs.lookup\_path(path) 递归解析路径 → 获取 inode → 查找或创建 File → fd\_manager.alloc(pid, file) → 返回 fd
    \item cat.asm：MOV R0,\#12 → INT 0x80（read）
    \item FileService：fd\_manager.get(fd) → OpenFile.read(count) → File.read(count) → 从 data[position..] 切片返回数据
    \item ProcessService：将返回数据写入进程虚拟内存 0x200，设 R2=字节数
    \item cat.asm：MOV R0,\#2 → INT 0x80（print\_str，从 0x200 打印 R2 字节）
\end{enumerate}

\textbf{VFS 持久化}：每次文件操作后，FileService 调用 vfs.save\_to\_file(".genshin-vfs.json") 将整个 VFS 树序列化为 JSON 写入磁盘。启动时调用 VirtualFileSystem::load\_from\_file 恢复。

\subsubsection{数据存储说明}

\begin{itemize}
    \item \texttt{.genshin-vfs.json} — VFS 树结构的 JSON 持久化文件，包含所有节点信息（inode、类型、父子关系、权限、磁盘块号）
    \item \texttt{.genshin-disk.img} — 虚拟磁盘镜像文件，512B 扇区，存储文件实际数据
\end{itemize}

\subsubsection{源程序文件说明}

\begin{itemize}
    \item \texttt{src/services/file/service.rs} — FileService 主逻辑
    \item \texttt{src/services/file/vfs.rs} — VFS 虚拟文件系统
    \item \texttt{src/services/file/file.rs} — File 文件抽象
    \item \texttt{src/services/file/descriptor.rs} — 文件描述符管理
    \item \texttt{src/hardware/disk.rs} — VirtualDisk 虚拟磁盘
    \item \texttt{src/hardware/block.rs} — 块设备层
\end{itemize}

\subsection{程序单元5：Shell 与用户界面}

\subsubsection{程序单元说明}

Shell 运行在主线程，提供交互式命令行界面。解析用户输入后通过消息总线发送请求，同步等待服务响应并打印结果。包含 pmon 进程监控器和 minigdb 调试器两个辅助工具。

\subsubsection{数据结构说明}

\begin{itemize}
    \item \textbf{Shell}（shell/mod.rs:38）：持有 UIContext（封装 MessageBus）、Timer 引用、当前工作目录 cwd
    \item \textbf{UIContext}（ui/mod.rs:18）：封装 MessageBus，提供 send(msg) 和 send\_request(msg) → Receiver<Response>
    \item \textbf{Command}（shell/parser.rs）：解析后的命令结构，包含 name 和 args
\end{itemize}

\subsubsection{算法及流程}

\textbf{Shell 命令执行流程}（以 ls 为例）：

\begin{enumerate}[label=(\arabic*)]
    \item 用户输入 \texttt{ls} → Parser 解析为 Command\{name: "ls", args: []\}
    \item execute\_command() → 匹配 "ls" → fork\_exec\_wait("ls", \&["/"])
    \item fork\_exec\_wait 内部：send\_and\_wait(ForkProcess) → 获取子进程 PID → send\_and\_wait(ExecProcess\{pid, "ls"\}) → send\_and\_wait(WaitChild\{pid, child\_pid\})
    \item ProcessService 处理 Fork → Exec（加载 ls.asm）→ 子进程运行 → MOV R0,\#18 → INT 0x80 → 结果写入 0x200 → MOV R0,\#2 → 打印结果 → HALT
    \item ProcessService 收到 WaitChild，等待子进程退出，回收资源
    \item Shell 收到响应，打印结果
\end{enumerate}

\subsubsection{源程序文件说明}

\begin{itemize}
    \item \texttt{src/ui/shell/mod.rs} — Shell 主逻辑（命令解析、执行、fork\_exec\_detach/wait、minigdb）
    \item \texttt{src/ui/shell/parser.rs} — 命令解析器
    \item \texttt{src/ui/shell/builtins.rs} — 内置命令
    \item \texttt{src/ui/shell/pmon.rs} — 进程监控器 TUI
    \item \texttt{src/ui/mod.rs} — UIContext
\end{itemize}

\subsection{程序单元6：硬件模拟层}

\subsubsection{程序单元说明}

硬件模拟层提供了操作系统运行所需的虚拟硬件：VirtualCPU 执行自定义 Mock ISA 指令、MMU 管理虚拟地址到物理地址的映射、Timer 以 100Hz 频率驱动进程调度、PhysicalMemory 和 VirtualDisk 提供存储。

\subsubsection{数据结构说明}

\begin{itemize}
    \item \textbf{VirtualCPU}（cpu.rs:260+）：registers[4]、pc、sp、flags（Zero/Sign/Carry/Overflow）、halted、syscall\_pending、syscall\_regs
    \item \textbf{MMU}（mmu.rs）：管理 per-PID 的页表，map\_page/unmap\_page/translate/get\_page\_entries
    \item \textbf{Timer}（timer.rs:51）：tick\_interval（10ms）、tick\_count、state（Stopped/Running/Paused），通过独立线程发送 Interrupt::Timer
    \item \textbf{PhysicalMemory}（memory.rs）：4MB 大数组模拟物理内存
\end{itemize}

\subsubsection{算法及流程}

\textbf{CPU 指令执行流程}（cpu.rs:371-404）：

\begin{enumerate}[label=(\arabic*)]
    \item fetch\_instruction()：从 MMU 读取 8 字节（opcode + dst + src\_type + pad + value）
    \item 解码 opcode（0x01=MOV, 0x03=SUB, 0x10=JMP, 0x13=JNZ, 0x80=INT, 0xFF=HALT 等）
    \item execute\_instruction()：执行指令，更新寄存器和标志位
    \item INT 0x80 特殊处理：保存寄存器到 syscall\_regs，设 syscall\_pending=true，不抛异常
\end{enumerate}

\textbf{定时器驱动调度}（timer.rs:107-142）：独立线程每 10ms 循环 → thread::sleep(10ms) → bus.send(Interrupt::Timer) → Kernel 路由到 ProcessService 的 intr\_rx → handle\_timer\_interrupt() 执行一轮调度。

\subsubsection{源程序文件说明}

\begin{itemize}
    \item \texttt{src/hardware/cpu.rs} — VirtualCPU（取指、解码、执行、中断处理）
    \item \texttt{src/hardware/mmu.rs} — MMU 地址转换
    \item \texttt{src/hardware/timer.rs} — Timer 定时器
    \item \texttt{src/hardware/memory.rs} — PhysicalMemory 物理内存
    \item \texttt{src/hardware/disk.rs} — VirtualDisk 虚拟磁盘
\end{itemize}

\newpage

% ═══════════════════════════════════════════════════════════
% 四、总结
% ═══════════════════════════════════════════════════════════

\section{总结}

genshin-OS 从零实现了一个微内核操作系统的核心组件。通过本项目，我们深入理解了以下操作系统核心概念：

\begin{enumerate}[label=(\arabic*)]
    \item \textbf{微内核架构}：内核只做消息路由，所有服务在用户态运行，通过消息总线通信。这种设计带来了良好的模块化和可扩展性，但也引入了消息传递的开销。
    \item \textbf{SMP 多核调度}：实现双 CPU 独立调度的过程中，遇到了 PID 重复分配、就绪队列竞态、量子过期处理等多个实际问题，加深了对调度器设计的理解。
    \item \textbf{同步原语与竞态}：CAS 无锁实现的信号量和互斥锁在高并发场景下表现出正确性。TOCTOU 所有权转移机制是解决多核环境下锁释放竞态的关键设计。
    \item \textbf{虚拟文件系统}：VFS 的 inode 树 + JSON 持久化方案在简单性和可用性之间取得了良好平衡。路径解析、文件描述符管理、磁盘块映射构成了完整的文件系统栈。
    \item \textbf{系统调用机制}：INT 0x80 → syscall\_pending → handle\_file\_syscall 分发的两层译码模型，在保持硬件层纯粹性的同时实现了高效的系统调用处理。
\end{enumerate}

项目开发过程中遇到的主要挑战包括：多核调度中进程被误判为 Zombie 导致提前回收（通过 fork 后不立即调度的策略修复）、信号量 count 在进程退出时被无条件 inflate 导致互斥失效（通过 count==0 前置检查和 holder 跟踪设计逐步改进）、消息总线的 SyscallTrap 泛滥导致服务无响应（通过外循环 Timer-only 过滤和内循环移除修复）等。这些问题的调试过程加深了我们对并发编程和操作系统内核设计复杂性的认识。

\newpage

\section{附录}

\subsection{源码文件清单}

\begin{center}
\begin{tabularx}{0.95\textwidth}{|l|X|}
\hline
\textbf{目录/文件} & \textbf{说明} \\
\hline
\texttt{src/main.rs} & 主入口，启动所有线程 \\
\texttt{src/messaging/} & 消息总线（bus.rs、msg.rs、response.rs） \\
\texttt{src/services/kernel.rs} & Kernel 路由器 \\
\texttt{src/services/process/} & 进程管理（service.rs、pcb.rs、scheduler.rs、ipc.rs、sync.rs、assembler.rs） \\
\texttt{src/services/memory/} & 内存管理（service.rs、alloc.rs、paging.rs、swap.rs） \\
\texttt{src/services/file/} & 文件系统（service.rs、vfs.rs、file.rs、descriptor.rs） \\
\texttt{src/services/device/} & 设备管理（service.rs、device.rs、driver.rs、manager.rs） \\
\texttt{src/hardware/} & 硬件模拟（cpu.rs、mmu.rs、timer.rs、memory.rs、disk.rs、block.rs、device.rs） \\
\texttt{src/ui/shell/} & Shell（mod.rs、parser.rs、builtins.rs、pmon.rs） \\
\texttt{src/ui/mod.rs} & UIContext \\
\texttt{programs/} & 用户态 .asm 测试程序 \\
\texttt{docs/} & 设计文档（ipc.md、syscalls.md、syncdemo.md 等） \\
\hline
\end{tabularx}
\end{center}

\subsection{参考文献}

\begin{enumerate}[label={[\arabic*]}]
    \item Rust 编程语言官方文档：\url{https://doc.rust-lang.org/book/}
    \item xv6 教学操作系统：\url{https://pdos.csail.mit.edu/6.828/2023/xv6.html}
    \item seL4 微内核：\url{https://sel4.systems/}
    \item 操作系统概念（第10版），Abraham Silberschatz 等
    \item The Design of the UNIX Operating System, Maurice J. Bach
\end{enumerate}

\end{document}
